\documentclass[11pt,oneside,a4paper,final]{article}

\usepackage{fancyhdr}
\usepackage{lastpage}

\pagestyle{fancy}

\renewcommand{\headrulewidth}{0.0pt}

\lhead{}
\chead{}
\rhead{}
\lfoot{}
\cfoot{Page {\thepage} of \pageref{LastPage}}
\rfoot{}

\usepackage{longtable}

\usepackage[style=numeric,backend=bibtex]{biblatex}
% \usepackage[style=numeric]{biblatex}
\bibliography{programmer_guide}
\usepackage{hyperref}
\usepackage[paper=a4paper,margin=1.6cm]{geometry}	% 1 inch margins all around
\usepackage[british=nohyphenation]{hyphsubst}
\usepackage[british]{babel}
\usepackage{graphicx}
\usepackage[caption=false,font=normalsize,labelfont=sf,textfont=sf]{subfig}
\usepackage{color}

\definecolor{mygreen}{rgb}{0,0.6,0}
\definecolor{mygray}{rgb}{0.5,0.5,0.5}
\definecolor{mymauve}{rgb}{0.58,0,0.82}

\usepackage{listings}
 \lstset{ %
  backgroundcolor=\color{white},   % choose the background color; you must add \usepackage{color} or \usepackage{xcolor}
  basicstyle=\footnotesize,        % the size of the fonts that are used for the code
  breakatwhitespace=false,         % sets if automatic breaks should only happen at whitespace
  breaklines=true,                 % sets automatic line breaking
  captionpos=b,                    % sets the caption-position to bottom
  commentstyle=\color{mygreen},    % comment style
%   deletekeywords={...},            % if you want to delete keywords from the given language
%   escapeinside={\%*}{*)},          % if you want to add LaTeX within your code
  extendedchars=true,              % lets you use non-ASCII characters; for 8-bits encodings only, does not work with UTF-8
  frame=single,                    % adds a frame around the code
  keepspaces=true,                 % keeps spaces in text, useful for keeping indentation of code (possibly needs columns=flexible)
  keywordstyle=\color{blue},       % keyword style
  language=C++,                 % the language of the code
%   morekeywords={*,...},            % if you want to add more keywords to the set
  numbers=left,                    % where to put the line-numbers; possible values are (none, left, right)
  numbersep=5pt,                   % how far the line-numbers are from the code
  numberstyle=\tiny\color{mygray}, % the style that is used for the line-numbers
  rulecolor=\color{black},         % if not set, the frame-color may be changed on line-breaks within not-black text (e.g. comments (green here))
  showspaces=false,                % show spaces everywhere adding particular underscores; it overrides 'showstringspaces'
  showstringspaces=false,          % underline spaces within strings only
  showtabs=false,                  % show tabs within strings adding particular underscores
  stepnumber=1,                    % the step between two line-numbers. If it's 1, each line will be numbered
  stringstyle=\color{mymauve},     % string literal style
  tabsize=2,                       % sets default tabsize to 2 spaces
  title=\lstname                   % show the filename of files included with \lstinputlisting; also try caption instead of title
}

\usepackage[acronym,toc]{glossaries}
\makeglossaries
% \phantomsection
% \addcontentsline{toc}{section}{Acronyms}
% \section*{Acronyms}
% \begin{footnotesize}
% \begin{acronym}[TDMA]

\newacronym{2D}{2D}{two-dimensional}
\newacronym{3D}{3D}{three-dimensional}
\newacronym{FLTK}{FLTK}{Fast Light Toolkit}
\newacronym{GLEW}{GLEW}{OpenGL Extension Wrangler Library}
\newacronym{GLSL}{GLSL}{OpenGL Shading Language}
\newacronym{GLUT}{GLUT}{OpenGL Utility Toolkit}
\newacronym{GPU}{GPU}{Graphics Processor Unit}
\newacronym{GUI}{GUI}{graphical user interface}
\newacronym{Gzip}{Gzip}{GNU zip}
\newacronym{ITK}{ITK}{Insight Segmentation and Registration Toolkit}
\newacronym{MXE}{MXE}{M cross environment}
\newacronym{SVN}{SVN}{Subversion}
\newacronym{VTK}{VTK}{Visualization Toolkit}



\title{gVirtualXRay -- Programmer's Guide}
\author{Dr Franck P. Vidal}
\date{18\textsuperscript{th} March 2014}

\begin{document}
 \sloppy

\maketitle

\newpage
\phantomsection
\addcontentsline{toc}{section}{Table of contents}
\tableofcontents

\newpage
\phantomsection
\addcontentsline{toc}{section}{List of figures}
\listoffigures

\newpage
\phantomsection
\addcontentsline{toc}{section}{List of tables}
\listoftables

\newpage
\phantomsection
\addcontentsline{toc}{section}{List of listings}
\lstlistoflistings

\newpage

\section{Introduction}

\section{Header inclusion}

\lstinputlisting[firstline=35, lastline=50, caption=Header inclusion., label=lst:header]{../../examples/src/tutorial.cxx}

Listing~\ref{lst:header} shows the header files that need to be included to build a simple program based on \Gls{GLUT}. 
\begin{itemize}
 \item \verb+glut.h+ is the main \Gls{GLUT} header file.
 \item \verb+iostream.h+ is used for output streams.
 \item \verb+exception.h+ is used for C++ exceptions.
 \item \verb+Types.h+ defines new types, e.g.~\verb+RATIONAL_NUMBER+, \verb+VEC2+, \verb+VEC3+ and \verb+MATRIX4+ to name a few.
 \item \verb+Units.h+ defines units such as metre, kilometre, electronvolt, kiloelectron volt, gram, kilogram, etc.
 \item \verb+TriangleMesh.h+ corresponds to a class used to handle \Acrshort{3D} triangle meshes.
 \item \verb+XRayBeam.h+ corresponds to a class used to manage X-Ray beams. The beam spectrum is discretised into energy channels.
 \item \verb+XRayDetector.h+ corresponds to a class that handles virtual X-Ray detectors.
 \item \verb+XRayRenderer.h+ is used to compute X-Ray images on the \Gls{GPU}.
\end{itemize}


\section{Name Spaces}

\lstinputlisting[firstline=55, lastline=57, caption=Name spaces., label=lst:name space]{../../examples/src/tutorial.cxx}

Listing~\ref{lst:name space} shows the name spaces that can be selected. 
\begin{itemize}
 \item \verb+Graphics+ is used to defined graphics elements such as \verb+TriangleMesh+. 
 \item \verb+gVirtualXRay+ is used to defined all the components of the X-ray simulation.
 \item \verb+std+ is used for output streams and exceptions.
\end{itemize}

\section{Setting Components of the X-Ray Simulation}

Figure~\ref{fig:scene} illustrates the main components of the simulation. 
\begin{figure}[tbh]
 \centering
 \includegraphics[width=0.5\textwidth]{xray-mesh}
 \caption{\label{fig:scene} Definition of the main components needed to simulated X-ray images.}
\end{figure}

Four of them have to be defined to simulate X-Ray images:
\begin{enumerate}
 \item The source (it includes the shape of the source, and the beam energy);
 \item The detector; and
 \item The scan objects.
\end{enumerate}


\subsection{X-Ray Detector}

\subsection{X-Ray Source}
Shape, Position, Beam,

\subsection{Scanned Object}

\section{Running the Simulation}

\section{Updating the Scanned Object}



% \newpage
% \section{Requirements}
% 
% The list of tools and libraries required to build gVirtualXRay is summarised in Table~\ref{tab:requiredToolsAndLibraries}.
% 
% \subsection{Main Library}
% 
% To build the library, you will need:
% \begin{itemize}
%  \item \href{http://www.cmake.org/}{\includegraphics[width=2cm]{500px-Cmake.png}}~\cite{cmake}\\
%  CMake is a cross-platform tool that we use to create the Unix Makefiles, XCode project files, MS VC++ project files, etc.~for the development platform that you may wish to use.
% 
%  \item \href{http://www.gzip.org/}{The \Gls{Gzip}}~\cite{gzip}\\
%  \Gls{Gzip} is used to compress the data required by the program at run-time (e.g. \Gls{GLSL} source code).
% 
%  \item \href{http://www.opengl.org/}{OpenGL}~\cite{opengl}\\
%  OpenGL is used to perform the simulation computations on the \Gls{GPU}.
% 
%  \item \href{http://glew.sourceforge.net/}{The \Gls{GLEW}}~\cite{glew}\\
%  \Gls{GLEW} is used to activate some of the OpenGL's extension that are needed by gVirtualXRay.
% 
%  \item \href{http://www.zlib.net/}{\includegraphics[width=2cm]{Zlib_3D_green}}~\cite{zlib}\\
%  The Zlib is used to decompress the data at run-time (the data compressed using \Gls{Gzip}).
% \end{itemize}
% 
% 
% \subsection{Examples and Demos}
% 
% To build the examples and demos, you will need:
% \begin{itemize}
%  \item \href{http://www.boost.org/}{\includegraphics[width=2cm]{Boost}}~\cite{boost}\\
%  It is only needed on the Mac OS X platform. 
%  It is mainly used to identify the path of the program when Bundles are used.
% 
%  \item An implementation of \href{http://www.opengl.org/resources/libraries/glut/}{The \Gls{GLUT}}~\cite{glut} (only \href{http://www.fltk.org/}{\includegraphics[width=2cm]{Fltk_shadow}}~\cite{fltk} or \href{http://freeglut.sourceforge.net/}{\includegraphics[width=2cm]{Freeglut_logo}}~\cite{freeglut} have been tested)\\
%  It is used to create the OpenGL's window of the demos.
% \end{itemize}
% 
% \begin{table}[tbhp]
% 	\begin{center}
% 		\caption{\label{tab:requiredToolsAndLibraries} Tools and libraries required to build gVirtualXRay.} 
% 		\begin{small}
% 			\begin{tabular}{|c|c|c|c|p{0.3\linewidth}|}
% 				\hline
% 				\textbf{Name} & \textbf{Component} & \textbf{Platform} & \textbf{Status} & \textbf{URL}\\
% 				\hline
% 				\hline
% 				CMake & Development tool & Linux, Mac OS X \& Windows & Required & \url{http://www.cmake.org/}\\
% 				\hline
% 				\Acrshort{Gzip} & System tool & Linux, Mac OS X \& Windows & Required & \url{http://www.gzip.org/}\\
% 				\hline
% 				OpenGL & Library & Linux, Mac OS X \& Windows & Required & \url{http://www.opengl.org/}\\
% 				\hline
% 				\Acrshort{GLEW} & Library & Linux, Mac OS X \& Windows & Required & \url{http://glew.sourceforge.net/}\\
% 				\hline
% 				Zlib & Library & Linux, Mac OS X \& Windows & Required & \url{http://www.zlib.net/}\\
% 				\hline
% 				Boost & Library & Mac OS X & Required & \url{http://www.boost.org/}\\
% 				\hline
% 				\Acrshort{GLUT} & Library & Linux, Mac OS X \& Windows & Required & \url{http://www.opengl.org/resources/libraries/glut/}\\
% 				\hline
% 				\Acrshort{FLTK} & Library & Linux, Mac OS X \& Windows & Optional & \url{http://www.fltk.org/}\\
% 				\hline
% 				freeglut & Library & Linux, Mac OS X \& Windows & Optional & \url{http://freeglut.sourceforge.net/}\\
% 				\hline
% 			\end{tabular}
% 		\end{small}
% 	\end{center}
% \end{table}
% 
% 
% \section{Options}
% 
% The list of optional tools and libraries to build gVirtualXRay is summarised in Table~\ref{tab:optionalToolsAndLibraries}.
% 
% \subsection{Main Library}
% 
% \begin{itemize}
%  \item \href{http://www.itk.org/}{\includegraphics[width=2cm]{InsightToolkitLogo}}~\cite{itk}\\
%  It is used to store \acrshort{2D} images or \acrshort{3D} volumes.
% 
%  \item\href{http://www.doxygen.org/}{\includegraphics[width=2cm]{Doxygen}}~\cite{doxygen}\\
% It is used to create the automatic documentation of the main library from the header files.
% \end{itemize}
% 
% \subsection{Examples and Demos}
% 
% 
% \begin{itemize}
%  \item \href{http://freeimage.sourceforge.net/}{FreeImage}~\cite{freeimage}\\
%  It is used to stored screenshots.
% \end{itemize}
% 
% 
% \begin{table}[tbhp]
% 	\begin{center}
% 		\caption{\label{tab:optionalToolsAndLibraries} Optional tools and libraries to build gVirtualXRay.} 
% 		\begin{small}
% 			\begin{tabular}{|c|c|c|c|p{0.3\linewidth}|}
% 				\hline
% 				\textbf{Name} & \textbf{Component} & \textbf{Platform} & \textbf{Status} & \textbf{URL}\\
% 				\hline
% 				\hline
% 				\Acrshort{ITK} & Library & Linux, Mac OS X \& Windows & Optional & \url{http://www.itk.org/}\\
% 				\hline
% 				Doxygen & Development tool & Linux, Mac OS X \& Windows & Optional & \url{http://www.doxygen.org/}\\
% 				\hline
% 				FreeImage & Library & Linux, Mac OS X \& Windows & Optional & \url{http://freeimage.sourceforge.net/}\\
% 				\hline
% 			\end{tabular}
% 		\end{small}
% 	\end{center}
% \end{table}
% 
% 
% \newpage
% \section{Compilation From the Source}
% 
% \subsection{Using \Acrshort{SVN}}
% Get the latest release from the \Gls{SVN} server:
% \begin{verbatim}
%    $ svn checkout svn://svn.code.sf.net/p/gvirtualxray/code/trunk gvirtualxray-code
% \end{verbatim}
% On Windows platforms, you may want to use \href{http://tortoisesvn.tigris.org/}{TortoiseSVN}~\cite{tortoisesvn}.
% 
% Set up the development environment using CMake.
% \begin{verbatim}
%    $ cd gvirtualxray-code
%    $ mkdir bin-release
%    $ cd bin-release
% \end{verbatim}
% 
% \subsection{Build the Project With Unix Makefiles}
% On Unix-like platforms (inc.~Linux and Mac OS X), to generate Makefiles, build the project, and install the library, type:
% \begin{verbatim}
%    $ cmake -G "Unix Makefiles" -DCMAKE_BUILD_TYPE=RELEASE ../build
%    $ make
%    $ make install  # You might need super-user RW rights
% \end{verbatim}
% 
% \subsection{Build the Project With XCode}
% On Mac OS X platforms, to generate XCode project, type:
% \begin{verbatim}
%    $ cmake -G "Xcode" -DCMAKE_BUILD_TYPE=RELEASE ../build
% \end{verbatim}
% 
% \subsection{Setting up the Project Using the \Acrshort{GUI}}
% On all the platforms (inc.~Windows), use the \gls{GUI} to have more options (e.g.~select different generators such as Visual C++, set variables). Then follow the instructions provided by the \Gls{GUI}:
% \begin{center}
% 	\begin{longtable}{|p{0.9\linewidth}|}
% 		\hline
% 		Select where the source (i.e.~CMakeLists.txt) is.
% 		Select where to build the binaries. This is where the Makefiles, XCode files or Visual C++ files will be created (depending on the generator that you will select).\\
% 		\includegraphics[width=0.5\linewidth]{cmake-01}\\
% 		\hline
% 		
% 		Click on \emph{Configure}:\\
% 		\includegraphics[width=0.5\linewidth]{cmake-02}\\
% 		\hline
% 
% 		Select the generator that you wish to use:\\
% 		\includegraphics[width=0.35\linewidth]{cmake-03}\\
% 		\hline
% 
% 		CMake will try to set up your development environment:\\
% 		\includegraphics[width=0.5\linewidth]{cmake-04}\\
% 		\hline
% 
% 		By default, the build type is Release, you can change it (with \emph{Debug}, \emph{MinSizeRel} or \emph{RelWithDebInfo}):\\
% 		\includegraphics[width=0.5\linewidth]{cmake-05}\\
% 		\hline
% 
% 		You can also manually modify or set the path to your libraries:\\
% 		\includegraphics[width=0.5\linewidth]{cmake-09}\\
% 		\hline
% 
% 		You can select different options, e.g. using \Gls{FLTK} or FreeGLUT (one or the other, not both), activate \Gls{ITK}:\\
% 		\includegraphics[width=0.5\linewidth]{cmake-06}\\
% 		\hline
% 
% 		After any change, you have to press \emph{Configure} again so that the changes are effective:\\
% 		\includegraphics[width=0.5\linewidth]{cmake-07}\\
% 		\hline
% 
% 		Once you have configured the project, click on \emph{Generate}.
% 		You are now ready to build the project with your selected development environment.\\
% 		\includegraphics[width=0.5\linewidth]{cmake-08}\\
% 		\hline
% 	\end{longtable}
% \end{center}
% 
% To make it a little easier on Windows platforms, it might be worth setting up environment variables to specify where you have installed some of the development tools/libraries. 
% Here is a list of the variables that may be useful to set:
% 
% \begin{itemize}
% 	\item \Gls{Gzip}:
% 	\begin{itemize}
% 		\item \verb|GZIP_TOOL|
% 	\end{itemize}
% 
% 	\item \Gls{GLEW}:
% 	\begin{itemize}
% 		\item \verb|GLEW_INCLUDE_DIR|
% 		\item \verb|GLEW_LIBRARY|
% 	\end{itemize}
% 	
% 	\item Zlib:
% 	\begin{itemize}
% 		\item \verb|ZLIB_INCLUDE_DIR|
% 		\item \verb|ZLIB_LIBRARY|
% 	\end{itemize}
% 	
% 	\item If you use FreeGLUT:
% 	\begin{itemize}
% 		\item \verb|GLUT_INCLUDE_DIR|
% 		\item \verb|GLUT_LIBRARY|
% 	\end{itemize}
% 
% 	\item Or if you use \Gls{FLTK}:
% 	\begin{itemize}
% 		\item \verb|FLTK_INCLUDE_DIR|
% 		\item \verb|FLTK_BASE_LIBRARY|
% 		\item \verb|FLTK_FLUID_EXECUTABLE|
% 		\item \verb|FLTK_FORMS_LIBRARY|
% 		\item \verb|FLTK_GL_LIBRARY|
% 		\item \verb|FLTK_IMAGES_LIBRARY|
% 		\item \verb|FLTK_MATH_LIBRARY|
% 	\end{itemize}
% 
% 	\item If you use \Gls{ITK}:
% 	\begin{itemize}
% 		\item \verb|ITK_DIR|
% 		\item And maybe \verb|VTK_DIR| if \Gls{ITK} was compiled with \Gls{VTK} support
% 	\end{itemize}
% 
% 	\item If you use FreeImage:
% 	\begin{itemize}
% 		\item \verb|FREEIMAGE_INCLUDE_DIR|
% 		\item \verb|FREEIMAGE_LIBRARY|
% 	\end{itemize}
% 
% \end{itemize}
% 
% 
% \section{Tested platforms}
% 
% Successful compilation and installation have been tested on:
% 
% \begin{itemize}
%  \item Linux openSUSE 12.3 (x86\_64) using the native compiler (g++ 4.7.2).
%  \item Linux openSUSE 12.3 (x86\_64) using the Windows cross compiler (mingw32-g++ 4.8.1) provided by \Gls{MXE}.
%  \item Windows 8.1 (x86\_64) using Visual C++ 2010 in 32 bits.
%  \item Windows 7 (x86) using Visual C++ 2013 in 32 bits.
%  \item Mac OS X 10.8 (x86\_64) using the native compiler (LLVM 4.2) in 64 bits.
%  \item Mac OS X 10.8 (x86\_64) using the GNU compiler (g++ 4.7) in 64 bits.
%  \item Mac OS X 10.8 (x86\_64) using the Windows cross compiler (mingw32-g++ 4.8.1) provided by \Gls{MXE}.
% \end{itemize}


\newpage
\printglossary[type=\acronymtype] 

\newpage
\phantomsection
\addcontentsline{toc}{section}{References}
\printbibliography


\end{document}
